\chapter{Introduction}
\label{chap:introduction}

Automotive radar is a rapidly evolving field of research, driven by the increasing demand for enhanced driver safety and the progressive realization of autonomous vehicles~\parencite{waldschmidtAutomotiveRadarFirst2021}.
While fully autonomous driving faces regulatory, economic, and technological challenges, the automotive industry has pragmatically prioritized vehicle safety via \glspl{adas}.
Among environmental perception modalities, millimetre-wave \gls{fmcw} radar operating in the \qtyrange{77}{81}{\giga\hertz} band remains uniquely resilient to adverse weather, extreme lighting, and low visibility.
However, while contemporary sensors excel at estimating target range, velocity, and angle, their target classification capability remains fundamentally limited.
Present commercial systems rely predominantly on single-polarized scalar \gls{mimo} arrays, optimized for geometric point-cloud generation but unequipped for detailed electromagnetic scattering characterization.
Consequently, distinguishing \glspl{vru} -- such as pedestrians, cyclists, and e-scooter riders -- from surrounding urban clutter and roadside infrastructure remains a critical bottleneck for high-level scene understanding~\parencite{tillyRoadUserClassification2021}.

In contrast to scalar systems, radar polarimetry captures the full vector nature of electromagnetic waves by measuring the transformation of polarization states upon target interaction.
Encoded within the complex $2 \times 2$ Sinclair scattering matrix $\vec S$, polarimetric responses reveal intrinsic structural, geometric, and material scatterer properties.
While target decomposition techniques have matured in satellite remote sensing and \gls{sar}, polarimetry has not yet been adopted in commercial automotive radar.
Yet, the transition to high-resolution 2D \gls{mimo} architectures at W-band offers a compelling opportunity: coupling polarimetric signatures with large virtual apertures and micro-Doppler spectra.
Resolving polarimetry across the Doppler domain enables the dynamic limb motion of pedestrians or the rotational dynamics of bicycle wheels to be isolated spectrally and characterized by their distinct physical scattering mechanisms~\parencite{bouwmeesterClassificationDynamicVulnerable2025}.

Despite its proven efficacy in remote sensing, radar polarimetry cannot be directly ingested into automotive radar without addressing severe domain transfer challenges.
The operational realities of dynamic driving environments -- grazing incidence propagation near the ground plane, low-altitude multipath, non-rigid target kinematics, mutual coupling, and phase-centre displacement across the array manifold -- systematically violate classical remote sensing assumptions.
A fundamental research gap exists: an absence of algorithm-driven hardware specification bounds, a lack of polarimetry-first \gls{mimo} topologies optimized for W-band, and a need for hybrid simulation frameworks and empirical campaigns capable of validating Doppler-resolved polarimetric models.

To resolve these challenges, this doctoral project establishes a top-down, algorithm-driven engineering methodology that bridges front-end antenna design, virtual array synthesis, and polarimetric signal processing.
The central goal of this research proposal is to investigate, formalize, and demonstrate polarimetric \gls{mimo} array architectures tailored for automotive radar.
To achieve this goal, the project is structured around five primary methodological pillars:
\begin{enumerate}
    \item \textit{Derivation of fundamental hardware bounds.} Execute perturbation analyses on decomposition algorithms to map breakdown points into minimum antenna specification bounds.
    \item \textit{Development of a hybrid simulation platform.} Construct an electromagnetic simulation framework linking full-wave antenna patterns with ray-tracing physics and radar signal synthesis.
    \item \textit{Synthesis of polarimetry-first MIMO topologies.} Formulate array optimization strategies that minimize phase-centre disparity and polarization squint across 2D virtual apertures.
    \item \textit{Design of high-purity W-band front-ends.} Develop planar and integrated antenna structures at \qty{79}{\giga\hertz} with high polarization purity and port isolation.
    \item \textit{Algorithm formulation and experimental validation.} Formulate Doppler-resolved feature extraction algorithms and validate target models through empirical measurement campaigns.
\end{enumerate}

\paragraph{Synopsis of the research proposal.}
The remainder of this research proposal is organized into five main chapters, mirroring the progression from theoretical foundations to operational execution.
\textbf{\Cref{chap:literature-review}} synthesizes the state of the art across radar polarimetry, high-performance W-band front-end antennas, and \gls{mimo} array multiplexing schemes.
\textbf{\Cref{chap:research-gap-analysis}} details the physical domain transfer challenges, categorizes identified literature gaps, and formalizes the four central research questions governing this doctoral work.
\textbf{\Cref{chap:feasibility-study}} outlines the proposed technical methodology while presenting preliminary feasibility results validating polarimetric ray tracing and W-band antenna array synthesis.
\textbf{\Cref{chap:research-planning}} establishes a four-year operational execution roadmap, mapping publishable topics to targeted venues alongside the formal doctoral education plan.
Finally, the proposal is concluded with a reflection on the candidate's research progress, methodology, and doctoral development.


